\begingroup
\renewcommand{\arraystretch}{1.45}
\SourceLeaf{312}
\Needspace{35\baselineskip}
\InterpolationArticle{25}
Wenn die Funktion $X$ nicht mit den Gliedern $\cos mx$, $\sin mx$ abbricht, sondern weiter fortläuft, und die Koeffizienten der folgenden Glieder noch zu beträchtlich sind, als dass man sie vernachlässigen dürfte, werden die Koeffizienten $\gamma,\gamma',\delta'$ usw. in der Funktion $X'$, die den $\mu$ Werten der Funktion $X$ genügt, welche für $x=a,b,c,d$ usw. eine vollständige Periode bilden, von den entsprechenden Koeffizienten $\alpha,\alpha',\beta'$ usw. merklich abweichen. Wenn wir also anstelle der Periode $x=a,b,c,d$ usw. eine andere Periode von $\mu$ Gliedern, etwa für $x=a',b',c',d'$ usw., auf ähnliche Weise behandeln und eine Funktion der $m^{\mathrm{ten}}$ Ordnung entwickeln, die diesen neuen Werten der Funktion $X$ genügt, und sie wie $X'$ durch den Ausdruck
\[
\begin{aligned}
&\gamma+\gamma'\cos x+\gamma''\cos2x+.\,.\,.+\gamma^m\cos mx\\[5pt]
&\phantom{\gamma}+\delta'\sin x+\delta''\sin2x+.\,.\,.+\delta^m\sin mx
\end{aligned}
\]
darstellen, so werden wir für die Koeffizienten $\gamma,\gamma',\delta'$ usw. nun Werte erhalten, die von den zuvor gefundenen merklich verschieden sind. Hieraus wird ersichtlich, wie diese Koeffizienten als \textit{veränderlich} betrachtet werden dürfen. Da nämlich aus den Artikeln 21, 23 hervorgeht, dass jeder von ihnen durch den Winkel $\mu a'$ ebenso bestimmt wird wie die früheren Werte durch den Winkel $\mu a$, kann jeder Koeffizient als Funktion des unbestimmten Winkels $y$ betrachtet werden, ähnlich wie $X$ eine Funktion des Winkels $x$ ist. \Needspace{34\baselineskip}
Nehmen wir nun an, dass $\mu\nu=\pi$ Werte der Funktion $X$, die eine vollständige Periode bilden, für
\[
\begin{aligned}
&x=a,a',a'',.\,.\,.b,b',b'',.\,.\,.c,c',c'',.\,.\,.d,d',d''.\,.\,.\text{ usw.}
\end{aligned}
\]
bekannt sind, sodass
\[
\begin{aligned}
&a'=a+\tfrac1\pi360^0.\\[5pt]
&a''=a+\tfrac2\pi360^0\text{ usw.}\\[5pt]
&b=a+\tfrac\nu\pi360^0=a+\tfrac1\mu360^0,\\[5pt]
&b'=a+\tfrac{\nu+1}\pi360^0=a'+\tfrac1\mu360^0\text{ usw.}\\[5pt]
&c=a+\tfrac{2\nu}\pi360^0=a+\tfrac2\mu360^0\text{ usw.}
\end{aligned}
\]
gilt und ersichtlich ist, dass sich diese Periode von $\pi$ Gliedern in $\nu$ Perioden von $\mu$ Gliedern für
\[
\begin{aligned}
&x=a,\ b,\ c,\ d\text{ usw.}\\[5pt]
&x=a',\ b',\ c',\ d'\text{ usw.}\\[5pt]
&x=a'',\ b'',\ c'',\ d''\text{ usw.}\\[5pt]
&\text{usw.}
\end{aligned}
\]
zerlegen lässt. Wenn wir also für jede einzelne Periode auf die gezeigte Weise die Funktion
\SourceLeaf{313}
$X'$ entwickeln, erhalten wir $\nu$ Werte jedes Koeffizienten $\gamma,\gamma',\delta'$ usw., die nacheinander den folgenden Werten von $y$ entsprechen:
\[
\begin{aligned}
&y=\mu a,\\[5pt]
& y=\mu a'=\mu a+\tfrac1\nu360^0,\\[5pt]
& y=\mu a''=\mu a+\tfrac2\nu360^0\text{ usw.}
\end{aligned}
\]
\textit{und daher vollständige Perioden von $\nu$ Gliedern bilden}. \Needspace{40\baselineskip}
Somit lässt sich nach unserer Methode jeder Koeffizient $\gamma,\gamma',\delta'$ usw. durch einen Ausdruck der Form
\[
\begin{aligned}
&\varepsilon+\varepsilon'\cos y+\varepsilon''\cos2y+.\,.\,.\,.+\varepsilon^n\cos ny\\[5pt]
&\phantom{\varepsilon}+\zeta'\sin y+\zeta''\sin2y+.\,.\,.\,.+\zeta^n\sin ny
\end{aligned}
\]
darstellen, wobei $n$ entweder $=\tfrac12\nu-\tfrac12$ oder $=\tfrac12\nu$ ist, je nachdem, ob $\nu$ ungerade oder gerade ist. Offenbar erhält ein solcher Ausdruck bei der Einsetzung von $y=\mu b,\mu c,\mu d$ usw. denselben Wert wie bei $y=\mu a$; ebenso bei $y=\mu b',\mu c',\mu d'$ usw. denselben wie bei $y=\mu a'$, und so weiter. Werden daher in
\[
\begin{aligned}
&\gamma+\gamma'\cos x+\gamma''\cos2x+\text{ usw.}\\[5pt]
&\phantom{\gamma}+\delta'\sin x+\delta''\sin2x+\text{ usw.}
\end{aligned}
\]
für die einzelnen Koeffizienten $\gamma,\gamma',\delta'$ usw. diese Funktionen des unbestimmten Winkels $y$ eingesetzt, so entsteht eine Funktion $Z$, die zwei Winkel $x,y$ enthält; dabei geht $Z$ in die aus der ersten Periode hergeleitete Funktion $X'$ über, wenn für $y$ einer der Werte $\mu a,\mu b,\mu c,\mu d$ usw. eingesetzt wird. Wird dagegen für $y$ einer der Werte $\mu a',\mu b',\mu c',\mu d'$ usw. genommen, so geht $Z$ in die aus der zweiten Periode gewonnene Funktion über. Hieraus folgt schließlich: Schreibt man in $Z$ für $y$ unmittelbar $\mu x$, sodass eine Funktion von $x$ allein entsteht ($X''$), so wird diese so beschaffen sein, dass sie für alle $\pi$ Werte von $x$ den vorgegebenen Werten der Funktion $X$ genügt.

\endgroup
